% Original PDF 22 / printed 12, complete Section 2.2.
% Fraktur source Z/B use the explicitly user-approved non-Fraktur mathcal rendering.
% block: cycles-boundaries-heading page:22 type:section
\subsection{{闭链, 边界链与同调}}\label{sec:cycles-boundaries-homology}

% block: cycles-boundaries-homology-definition page:22 type:ordinary-paragraph
{给定链复形 $(V,\partial)$, $\partial$ 的零空间 $\mathcal Z$ 和像空间 $\mathcal B$ 都是 $V$ 的分次子空间. (它们也是子复形, 不过其微分相当乏味, 因为 $\partial$ 在两者上的限制均为零.) $\mathcal Z_k$ 的元素称为 $k$-闭链, 而 $\partial_{k+1}$ 的像空间 $\mathcal B_k$ 的元素称为 $k$-边界链. 微分的性质 $\partial^2=0$ 意味着 $\mathcal B_k\subset\mathcal Z_k$, 第 $k$ 个\emph{同调空间}定义为商空间 $\mathcal H_k=\mathcal Z_k/\mathcal B_k$. 因此, 同调空间的元素是闭链的等价类: 若两个 $k$-闭链之差是某个 $(k+1)$-链的边界, 就认为它们等价. 若边界链空间 $\mathcal B_k$ 恰好等于闭链空间 $\mathcal Z_k$, 而非其真子空间, 则所有同调空间 $\mathcal H_k$ 都为零. 此时称该复形是\emph{正合的}.}
